\uuid{SU23}
\exo7id{7077}
\titre{exo7 7077}
\auteur{megy}
\organisation{exo7}
\datecreate{2017-01-21}
\isIndication{false}
\isCorrection{false}
\chapitre{Géométrie affine euclidienne}
\sousChapitre{Géométrie affine euclidienne du plan}
\cpge{mpsi-mp2i.applications-lineaires.b, pcsi.applications-lineaires.b, ptsi.applications-lineaires.b, mpsi-mp2i.espaces-prehilbertiens.e, pcsi.espaces-prehilbertiens.e, bcpst1.geometrie.e, mp-mpi.endomorphismes-euclidien.c, pc.endomorphismes-euclidiens.a, psi.endomorphismes-euclidien.a, pt.isometries-euclidien.a, mp-mpi.endomorphismes-euclidien.d, pc.endomorphismes-euclidiens.c, psi.endomorphismes-euclidien.d, pt.isometries-euclidien.c}
\module{Géométrie}
\niveau{L2}
\difficulte{}

\contenu{
\question{
% Source :
% Tags :  
Soit $\mathcal D$ une droite et $\sigma$ la réflexion orthogonale suivant cette droite. Construire l'image par $\sigma$ d'un point, d'un segment, d'un droite, d'un cercle (que ces objets intersectent l'axe  de symétrie ou pas).
}
}
